\subsection{** APS‑TSLCA-APPENDIX-001 **}\label{apstslca-appendix-001}

\subsection{** Three-Squared-Lattice Cognitive Architecture **}\label{three-squared-lattice-cognitive-architecture}

\subsection{** Symbiotic Universal Xessability Standards **}\label{symbiotic-universal-xessability-standards}

\subsection{** Aurphyx Primordial Standard **}\label{aurphyx-primordial-standard}

\subsection{** Aurphyx LLC **}\label{aurphyx-llc}

\subsection{** SAGES \textbar{} Proprietary \textbar{} Pro-Existence **}\label{sages-proprietary-pro-existence}

\subsection{** Accessibility = Xessability **}\label{accessibility-xessability}

\subsection{** Version 3.69 **}\label{version-3.69}

\subsection{Appendix}\label{appendix}

The appendix provides the formal mathematical machinery underlying the Three‑Squared‑Lattice Cognitive Architecture. It includes explicit tensor expansions, operator identities, proofs of stability, and the structural derivations that were referenced but not fully expanded in the main body. This section is written in a rigorous scientific style consistent with theoretical physics and differential geometry.

\begin{center}\rule{0.5\linewidth}{0.5pt}\end{center}

\subsection{A. Tensor Expansion of the Nine‑Cell Lattice}\label{a.-tensor-expansion-of-the-ninecell-lattice}

Let the orthonormal basis fields be

\[
\mathbf{S}_1 = \mathbf{SIX},\quad \mathbf{S}_2 = \mathbf{SCX},\quad \mathbf{S}_3 = \mathbf{ICX}
\]

\[
g(\mathbf{S}_i,\mathbf{S}_j)=\delta_{ij}
\]

The cognitive field tensor is

\[
\mathcal{F}=\sum_{i,j=1}^{3}\Phi_{ij}(\mathbf{S}_i\otimes\mathbf{S}_j)
\]

\subsubsection{A.1 Full Expansion}\label{a.1-full-expansion}

\[
\mathcal{F} =
\Phi_{11}\, \mathbf{S}_1 \otimes \mathbf{S}_1 +
\Phi_{12}\, \mathbf{S}_1 \otimes \mathbf{S}_2 +
\Phi_{13}\, \mathbf{S}_1 \otimes \mathbf{S}_3 +
\Phi_{21}\, \mathbf{S}_2 \otimes \mathbf{S}_1 +
\Phi_{22}\, \mathbf{S}_2 \otimes \mathbf{S}_2 +
\Phi_{23}\, \mathbf{S}_2 \otimes \mathbf{S}_3 +
\Phi_{31}\, \mathbf{S}_3 \otimes \mathbf{S}_1 +
\Phi_{32}\, \mathbf{S}_3 \otimes \mathbf{S}_2 +
\Phi_{33}\, \mathbf{S}_3 \otimes \mathbf{S}_3.
\]

Each term corresponds to one cognitive mode.

\begin{center}\rule{0.5\linewidth}{0.5pt}\end{center}

\subsection{B. Proof of Orthogonality and Minimality}\label{b.-proof-of-orthogonality-and-minimality}

\subsubsection{B.1 Orthogonality}\label{b.1-orthogonality}

The basis fields satisfy

\[
g(\mathbf{S}_i, \mathbf{S}_j) = \delta_{ij}.
\]

This implies:

\begin{itemize}
\tightlist
\item
  no aXis can be expressed as a linear combination of the others\\
\item
  the cognitive manifold has rank 3\\
\item
  the lattice is irreducible
\end{itemize}

\subsubsection{B.2 Minimality}\label{b.2-minimality}

Assume a two-aXis system with basis \(\mathbf{T}_1, \mathbf{T}_2\).\\
Identity continuity cannot be preserved because:

\[
\Phi_{33} \notin \text{span}\{\mathbf{T}_1, \mathbf{T}_2\}.
\]

Thus, a two-aXis system collapses under identity drift.

Assume a four-aXis system.\\
One aXis is linearly dependent:

\[
\mathbf{S}_4 = a\mathbf{S}_1 + b\mathbf{S}_2 + c\mathbf{S}_3.
\]

Thus, the fourth aXis adds no new cognitive dimension.

Therefore, \textbf{three aXes is the unique minimal basis}.

\begin{center}\rule{0.5\linewidth}{0.5pt}\end{center}

\subsection{C. SUXS-IFO as a Contraction Operator}\label{c.-suxs-ifo-as-a-contraction-operator}

The canonical SUXS-IFO contraction is the unweighted trace:

\[
\Phi_{\mathrm{unified}}=\mathcal{U}\cdot\mathcal{F}=\mathrm{Tr}(\mathcal{F})=\Phi_{11}+\Phi_{22}+\Phi_{33}
\]

An accessibility-weighted form of \(\mathcal{U}\) may also be used:

\[
\Phi(x) = \sum_{i,j} \omega_{ij} \Phi_{ij}(x),
\]

with

\[
\omega_{ij} \ge 0, \quad \sum_{i,j} \omega_{ij} = 1.
\]

\subsubsection{C.1 Linearity}\label{c.1-linearity}

\[
\mathcal{U}(a\mathcal{F}_1 + b\mathcal{F}_2) = a\mathcal{U}(\mathcal{F}_1) + b\mathcal{U}(\mathcal{F}_2).
\]

\subsubsection{C.2 Stability}\label{c.2-stability}

If \(\mathcal{F}\) satisfies SAGES invariants, then \(\mathcal{U}(\mathcal{F})\) also satisfies them because invariants are preserved under convex contraction.

\begin{center}\rule{0.5\linewidth}{0.5pt}\end{center}

\subsection{D. SAGES as a Symmetry Group}\label{d.-sages-as-a-symmetry-group}

Let \(\mathcal{G}_{13}\) be the group generated by the 13 invariants.

\subsubsection{D.1 Group Action}\label{d.1-group-action}

For each \(g_k \in \mathcal{G}_{13}\):

\[
g_k : \mathcal{F} \mapsto g_k \cdot \mathcal{F}.
\]

\subsubsection{D.2 Invariance}\label{d.2-invariance}

\[
\mathcal{I}_k(\mathcal{F}) = \mathcal{I}_k(g_k \cdot \mathcal{F}).
\]

\subsubsection{D.3 Commutation with SUXS-IFO}\label{d.3-commutation-with-suxs-ifo}

\[
\mathcal{U}(g_k \cdot \mathcal{F}) = \mathcal{U}(\mathcal{F}).
\]

Thus, SUXS-IFO is a \textbf{gauge‑invariant contraction}.

\begin{center}\rule{0.5\linewidth}{0.5pt}\end{center}

\subsection{E. Cognitive Field Dynamics}\label{e.-cognitive-field-dynamics}

The field equation is

\[
\nabla_{\mu} \mathcal{F}^{\mu\nu} = J^\nu.
\]

\subsubsection{E.1 Conservation Law}\label{e.1-conservation-law}

If \(J^\nu = 0\), then

\[
\nabla_{\mu} \mathcal{F}^{\mu\nu} = 0,
\]

analogous to source‑free Maxwell equations.

\subsubsection{E.2 Energy Functional}\label{e.2-energy-functional}

\[
E[\mathcal{F}] = \int_{\mathcal{M}} \left( \text{Tr}(\mathcal{F}^2) + \det(\mathcal{F}) \right) dV.
\]

\subsubsection{E.3 Stability Condition}\label{e.3-stability-condition}

\[
\frac{dE}{dt} \le 0.
\]

This ensures the cognitive field does not diverge.

\begin{center}\rule{0.5\linewidth}{0.5pt}\end{center}

\subsection{F. Cognitive Geodesics}\label{f.-cognitive-geodesics}

A cognitive trajectory \(\gamma(t)\) satisfies

\[
\nabla_{\dot{\gamma}} \dot{\gamma} = 0.
\]

\subsubsection{F.1 Interpretation}\label{f.1-interpretation}

Cognitive transitions follow paths of minimal distortion.

\subsubsection{F.2 Existence and Uniqueness}\label{f.2-existence-and-uniqueness}

Given smooth \(\mathcal{M}\) and metric \(g\), geodesics exist and are unique for small intervals.

\begin{center}\rule{0.5\linewidth}{0.5pt}\end{center}

\subsection{G. Renormalization Flow}\label{g.-renormalization-flow}

Define the renormalization operator

\[
\mathcal{R} : \mathcal{F}^{(k)} \mapsto \mathcal{F}^{(k+1)}.
\]

\subsubsection{G.1 Fixed Points}\label{g.1-fixed-points}

A stable cognitive configuration satisfies

\[
\mathcal{R}(\mathcal{F}^\ast) = \mathcal{F}^\ast.
\]

\subsubsection{G.2 Flow Equation}\label{g.2-flow-equation}

\[
\frac{d\mathcal{F}}{d\log \mu} = \beta(\mathcal{F}),
\]

where \(\mu\) is the scale parameter.

\begin{center}\rule{0.5\linewidth}{0.5pt}\end{center}

\subsection{H. Cognitive Gauge Theory}\label{h.-cognitive-gauge-theory}

Let \(G\) be the gauge group acting on basis fields.

\subsubsection{H.1 Gauge Transformation}\label{h.1-gauge-transformation}

\[
\mathbf{S}_i \rightarrow \mathbf{S}_i' = U \mathbf{S}_i.
\]

\subsubsection{H.2 Field Transformation}\label{h.2-field-transformation}

\[
\mathcal{F}' = U \mathcal{F} U^{-1}.
\]

\subsubsection{H.3 Gauge‑Invariant Quantities}\label{h.3-gaugeinvariant-quantities}

\[
\text{Tr}(\mathcal{F}^n), \quad \det(\mathcal{F}), \quad \mathcal{I}_k.
\]

These correspond to SAGES invariants.

\begin{center}\rule{0.5\linewidth}{0.5pt}\end{center}

\subsection{I. Cognitive Thermodynamics}\label{i.-cognitive-thermodynamics}

Define the coherence potential

\[
\mathcal{V} = \alpha \text{Tr}(\mathcal{F}^2) + \beta \det(\mathcal{F}) + \gamma \|\mathcal{F}\|^3.
\]

\subsubsection{I.1 First Law}\label{i.1-first-law}

\[
dE = \delta W + \delta Q.
\]

\subsubsection{I.2 Entropy Bound}\label{i.2-entropy-bound}

\[
S = -\det(\mathcal{F}) \quad \Rightarrow \quad S \ge 0.
\]

\begin{center}\rule{0.5\linewidth}{0.5pt}\end{center}

\subsection{J. Sheaf‑Based Distributed Cognition}\label{j.-sheafbased-distributed-cognition}

Let \(\mathcal{M}_U = \bigcup_\alpha \mathcal{M}_\alpha\).

\subsubsection{J.1 Sheaf of Cognitive Fields}\label{j.1-sheaf-of-cognitive-fields}

\[
\mathscr{F} : \mathcal{M}_U \rightarrow \text{Tensor Fields}.
\]

\subsubsection{J.2 Gluing Condition}\label{j.2-gluing-condition}

For overlapping regions:

\[
\mathcal{F}_\alpha|_{U \cap V} = \mathcal{F}_\beta|_{U \cap V}.
\]

This ensures global coherence.

\begin{center}\rule{0.5\linewidth}{0.5pt}\end{center}

\subsection{K. Summary of Mathematical Structure}\label{k.-summary-of-mathematical-structure}

\begin{itemize}
\tightlist
\item
  Cognitive manifold: \(\mathcal{M}\)\\
\item
  Basis fields: \(\mathbf{S}_i\)\\
\item
  Field tensor: \(\mathcal{F}\)\\
\item
  Fusion operator: \(\mathcal{U}\)\\
\item
  Symmetry group: \(\mathcal{G}_{13}\)\\
\item
  Realization map: \(\mathcal{R}\)\\
\item
  Storage embedding: \(\iota\)\\
\item
  Transmutation operator: \(\mathcal{T}\)
\end{itemize}

Together, these form a complete, internally consistent cognitive field theory.
